\documentclass[12pt]{article}
\usepackage[T1]{fontenc}
\usepackage[utf8]{inputenc}
\usepackage{lmodern}
\usepackage{textcomp}
\usepackage{enumitem}
\usepackage{geometry}
\geometry{margin=1in}
\begin{document}
\begin{center}
Answer Key - Mathematics - Undergraduate Level Assessment 1 \
\small (Answers are ordered exactly as in the question paper.)
\end{center}

\section*{Multiple Choice}
\begin{enumerate}[label=\arabic*.]
\item \textbf{Answer:} B
\\ \textit{Options:} A) Yes, with p=2.; B) Yes, with p=3.; C) Yes, with p=5.; D) No.
\\ \textit{Note:} Correct option: B - Yes, with p=3.
\item \textbf{Answer:} D
\\ \textit{Options:} A) True, True; B) False, False; C) True, False; D) False, True
\\ \textit{Note:} Correct option: D - False, True
\item \textbf{Answer:} B
\\ \textit{Options:} A) If A is true, then B is false.; B) If A is false, then B is false.; C) If A is false, then B is true.; D) Both A and B are true.
\\ \textit{Note:} Correct option: B - If A is false, then B is false.
\item \textbf{Answer:} B
\\ \textit{Options:} A) 24; B) 30; C) 32; D) 36
\\ \textit{Note:} Correct option: B - 30
\item \textbf{Answer:} A
\\ \textit{Options:} A) 0; B) 3; C) 12; D) 30
\\ \textit{Note:} Correct option: A - 0
\end{enumerate}

\section*{True / False}
\begin{enumerate}[label=\arabic*.]
\setcounter{enumi}{5}
\item \textbf{Answer:} False
\\ \textit{Options:} A) True; B) False
\\ \textit{Note:} LLM-generated false statement. Correct answer: '-1'.
\item \textbf{Answer:} True
\\ \textit{Options:} A) True; B) False
\\ \textit{Note:} LLM-generated true statement based on MCQ.
\item \textbf{Answer:} False
\\ \textit{Options:} A) True; B) False
\\ \textit{Note:} LLM-generated false statement. Correct answer: '$e^2$/2'.
\item \textbf{Answer:} True
\\ \textit{Options:} A) True; B) False
\\ \textit{Note:} LLM-generated true statement based on MCQ.
\item \textbf{Answer:} False
\\ \textit{Options:} A) True; B) False
\\ \textit{Note:} LLM-generated false statement. Correct answer: 'a point'.
\end{enumerate}

\section*{Long Form Response}
\begin{enumerate}[label=\arabic*.]
\setcounter{enumi}{10}
\item \textbf{Answer:} The remainder when $37+49+14$ is divided by $12$ is the same as the remainder when the congruences of each number mod 12 are summed and then divided by 12. In other words, since we have \begin{align*} 37 \&\ensuremath{\equiv} 1\ensuremath{\pmod{12}}\\ 49 \&\ensuremath{\equiv} 1\ensuremath{\pmod{12}} \\ 14 \&\ensuremath{\equiv} 2\ensuremath{\pmod{12}} \end{align*}then therefore, $37+49+14 \equiv 1+1+2 \equiv \boxed{4}\pmod{12}$.
\\ \textit{Note:} Students should show all steps in their mathematical solution.
\item \textbf{Answer:} 35. Explain why this is the correct answer and provide supporting evidence. Show the calculation or reasoning that leads to this conclusion.
\\ \textit{Note:} Long-form derived from MCQ; award credit for stating the correct idea and giving supporting reasoning.
\end{enumerate}

\vfill
\noindent\textit{End of Answer Key}
\end{document}